\documentclass[11pt]{article}
\usepackage[a4paper,margin=1in]{geometry}
\usepackage{fontspec}
\usepackage[boldfont=false]{xeCJK}
\setCJKmainfont{FandolSong-Regular.otf}
\usepackage{amsmath}
\usepackage{amssymb}
\usepackage{array}
\usepackage{booktabs}
\usepackage{graphicx}
\usepackage{adjustbox}
\usepackage{enumitem}
\setlist[itemize]{topsep=2pt,partopsep=0pt,parsep=0pt,itemsep=2pt,leftmargin=1.4em}
\setlist[enumerate]{topsep=2pt,partopsep=0pt,parsep=0pt,itemsep=2pt,leftmargin=1.6em}
\usepackage{parskip}
\usepackage{fvextra}
\fvset{breaklines=true,breakanywhere=true,fontsize=\small}
\setlength{\parindent}{0pt}

\begin{document}

\begin{center}
  {\LARGE\bfseries Human nutrition}\\[0.35em]
  {\large Igcse Biology}
\end{center}
\vspace{0.6em}

\section*{A balanced diet}

A \textbf{balanced diet} 均衡饮食 gives you all the substances you need, in the right amounts. There are seven parts.

\begin{center}
\adjustbox{max width=\linewidth}{%
\begin{tabular}{lll}
\toprule
\textbf{Part of diet} & \textbf{Good sources} & \textbf{Why you need it} \\
\midrule
\textbf{carbohydrates} 碳水化合物 & rice, bread, potatoes & the main source of \textbf{energy} 能量 \\
\textbf{fats} 脂肪 and oils & butter, oil, nuts & an energy store; help keep you warm \\
\textbf{proteins} 蛋白质 & meat, fish, eggs, beans & growth and repair of cells \\
\textbf{vitamin} 维生素 C & fresh fruit, vegetables & keeps skin and gums healthy \\
vitamin D & sunlight, oily fish, eggs & helps the body take in \textbf{calcium} 钙 \\
\textbf{mineral ions} 矿物质离子 — calcium & milk, cheese & strong bones and teeth \\
mineral ions — \textbf{iron} 铁 & red meat, leafy greens & needed to make \textbf{red blood cells} 红细胞 that carry oxygen \\
\textbf{fibre} 膳食纤维 (roughage) & vegetables, whole grains & keeps food moving through the gut \\
water & drinks and food & needed for all reactions; most of the body is water \\
\bottomrule
\end{tabular}}
\end{center}

Two diseases come from missing a vitamin:

\begin{itemize}
\item too little vitamin C → \textbf{scurvy} 坏血病 (bleeding gums, slow healing).

\item too little vitamin D → \textbf{rickets} 佝偻病 (soft, bent bones), because the body cannot take in enough calcium.

\end{itemize}

\section*{The digestive system}

The \textbf{digestive system} 消化系统 breaks food down. Food passes along one long tube, the \textbf{alimentary canal} 消化道:

mouth → \textbf{oesophagus} 食道 → \textbf{stomach} 胃 → \textbf{small intestine} 小肠 (the \textbf{duodenum} 十二指肠 then the \textbf{ileum} 回肠) → \textbf{large intestine} 大肠 (the \textbf{colon} 结肠, \textbf{rectum} 直肠 and \textbf{anus} 肛门).

Some organs help with digestion but food does not pass through them: the \textbf{salivary glands} 唾液腺 in the mouth, the \textbf{pancreas} 胰腺, the \textbf{liver} 肝脏 and the \textbf{gall bladder} 胆囊.

Five things happen to food, in order:

\begin{center}
\adjustbox{max width=\linewidth}{%
\begin{tabular}{ll}
\toprule
\textbf{Process} & \textbf{What it means} \\
\midrule
\textbf{ingestion} 摄入 & taking food and drink into the body through the mouth \\
\textbf{digestion} 消化 & breaking food down into small molecules \\
\textbf{absorption} 吸收 & \textbf{nutrients} 营养物质 move from the intestine into the blood \\
\textbf{assimilation} 同化 & cells take in and use the nutrients \\
\textbf{egestion} 排遗 & undigested food leaves the body as \textbf{faeces} 粪便 \\
\bottomrule
\end{tabular}}
\end{center}

\section*{Physical digestion}

Physical digestion breaks food into smaller pieces \textbf{without} changing the food molecules. This gives a larger \textbf{surface area} 表面积 for the \textbf{enzymes} 酶 to act on later.

\subsection*{Teeth}

You have four kinds of \textbf{teeth} 牙齿:

\begin{center}
\adjustbox{max width=\linewidth}{%
\begin{tabular}{ll}
\toprule
\textbf{Tooth} & \textbf{Job} \\
\midrule
\textbf{incisors} 门齿 & sharp front teeth for biting and cutting \\
\textbf{canines} 犬齿 & pointed teeth for tearing \\
\textbf{premolars} 前臼齿 & flat teeth for chewing and grinding \\
\textbf{molars} 臼齿 & flat back teeth for chewing and grinding \\
\bottomrule
\end{tabular}}
\end{center}

A tooth is built from these parts:

\begin{center}
\adjustbox{max width=\linewidth}{%
\begin{tabular}{ll}
\toprule
\textbf{Part} & \textbf{Description} \\
\midrule
\textbf{enamel} 牙釉质 & hard white outer layer; the hardest material in the body \\
\textbf{dentine} 牙本质 & softer, bone-like layer under the enamel \\
\textbf{pulp} 牙髓 & soft centre with \textbf{nerves} 神经 and \textbf{blood vessels} 血管 \\
\textbf{cement} 牙骨质 & fixes the root into the jaw \\
\bottomrule
\end{tabular}}
\end{center}

Teeth are set into the bone of the jaw and held firm by the \textbf{gums} 牙龈.

\subsection*{The stomach}

The wall of the stomach is made of \textbf{muscle} 肌肉 that squeezes and mixes the food, breaking it into smaller pieces.

\subsection*{Bile (Supplement)}

\textbf{Bile} 胆汁 is made in the liver and stored in the gall bladder. It \textbf{emulsifies} 乳化 fats and oils — it breaks large drops of fat into many tiny droplets. This gives a much larger surface area for the enzyme \textbf{lipase} 脂肪酶 to digest the fat. Bile itself contains no enzymes.

\section*{Chemical digestion}

Chemical digestion breaks large \textbf{insoluble} 不溶性 molecules into small \textbf{soluble} 可溶性 molecules. Only small soluble molecules can be absorbed into the blood. Enzymes carry out this work.

\begin{center}
\adjustbox{max width=\linewidth}{%
\begin{tabular}{llll}
\toprule
\textbf{Enzyme} & \textbf{Breaks down} & \textbf{Into} & \textbf{Made in} \\
\midrule
\textbf{amylase} 淀粉酶 & \textbf{starch} 淀粉 & \textbf{reducing sugars} 还原糖 & salivary glands, pancreas \\
\textbf{protease} 蛋白酶 & protein & \textbf{amino acids} 氨基酸 & stomach, pancreas \\
lipase & fats and oils & \textbf{fatty acids} 脂肪酸 and \textbf{glycerol} 甘油 & pancreas \\
\bottomrule
\end{tabular}}
\end{center}

\subsection*{Acid in the stomach}

The stomach makes \textbf{hydrochloric acid} 盐酸, which is part of the \textbf{gastric juice} 胃液. This acid:

\begin{itemize}
\item kills harmful \textbf{microorganisms} 微生物 in the food, and

\item gives an \textbf{acidic} 酸性 pH, the best pH for the stomach's protease to work.

\end{itemize}

\subsection*{Digesting starch and protein (Supplement)}

Starch is digested in two steps:

\begin{itemize}
\item amylase breaks starch into \textbf{maltose} 麦芽糖.

\item \textbf{maltase} 麦芽糖酶, on the membranes of the \textbf{epithelium} 上皮 lining the small intestine, breaks maltose into \textbf{glucose} 葡萄糖.

\end{itemize}

Protein is digested by two proteases:

\begin{itemize}
\item \textbf{pepsin} 胃蛋白酶 breaks protein down in the acidic stomach.

\item \textbf{trypsin} 胰蛋白酶 breaks protein down in the \textbf{alkaline} 碱性 small intestine.

\end{itemize}

Bile is alkaline. It \textbf{neutralises} 中和 the acidic food and gastric juice as they enter the duodenum, giving a suitable pH for the enzymes there.

\section*{Absorption}

Digested nutrients are absorbed into the blood in the small intestine. Most water is absorbed here too; the colon absorbs the rest, leaving solid faeces.

\subsection*{Villi (Supplement)}

The inside of the small intestine is covered with millions of tiny finger-shaped \textbf{villi} 绒毛. Each villus is itself covered in much smaller \textbf{microvilli} 微绒毛. Together these give a huge surface area, so absorption is fast.

A single villus has:

\begin{itemize}
\item a wall just one cell thick, so molecules have only a short distance to travel.

\item a network of \textbf{capillaries} 毛细血管 that carry away absorbed glucose and amino acids.

\item a \textbf{lacteal} 乳糜管 in the centre that carries away absorbed fatty acids and glycerol.

\end{itemize}

\section*{Exam tips}

\begin{itemize}
\item A balanced diet has the right amount of each part. Learn one source and one use for each.

\item Vitamin C ↔ scurvy; vitamin D ↔ rickets.

\item The five processes in order: ingestion → digestion → absorption → assimilation → egestion. Egestion (faeces) is \textbf{not} excretion.

\item Physical digestion makes pieces smaller (more surface area); chemical digestion uses enzymes to make molecules smaller and soluble.

\item amylase → sugars, protease → amino acids, lipase → fatty acids + glycerol. Stomach acid kills microorganisms and sets a low pH.

\item Villi and microvilli give a large surface area, and a thin wall gives a short distance — both make absorption faster.

\end{itemize}


\end{document}
